\section{The geometric systems and a finite list of correspondences}
\label{geo:section}

We first fix the curves, the family, and an elliptic-square center.  The
trace-free self-adjoint endomorphisms form a five-dimensional orthogonal
system.  Its monodromy gives an alternative for two varying points: the
two systems have product monodromy, or their joint image lies in a
polarized isogeny correspondence.  The final result of this section
reduces the latter alternative to a fixed finite list.  This list is
chosen before the auxiliary levels of \Cref{inc:markings}.

\subsection{Fixed covers and the center}

\begin{proposition}\label{geo:setup}
In proving \Cref{intro:squares}, we may work with a smooth geometrically
irreducible curve $S$ over a fixed number field, a principally polarized
abelian scheme $\mathscr A\to S$, and a nonconstant finite map from $S$
onto a dense open of the original curve $C$.  We may equip the family
with full symplectic levels at two distinct fixed primes at least three.
Under the supposition that the required elliptic-square locus is
infinite, we may fix a point $t_0\in S(\Qbar)$ in that locus.

Further fixed finite covers, normalization, and removal of finite sets
preserve this reduction.  After enlarging a fixed field of definition
$K$, there is a constant $b$ such that every selected lift $t$ of a
retained coarse point $s$ satisfies
\begin{equation}\label{geo:degree-descent}
 [K(s):K]\leq [K(t):K]\leq b[K(s):K].
\end{equation}
The smooth projective completion of $S$ and the divisor $[t_0]$ are
available even when $C$ is complete.
\end{proposition}

\begin{proof}
Choose distinct primes $r_1,r_2\geq3$ and put $N=r_1r_2$.
These preliminary levels give a fine moduli cover.  The auxiliary
incidence primes will be chosen separately in \Cref{inc:markings},
after the correspondence list has been fixed.
After adjoining the requisite roots of unity, the fine Siegel moduli
space $\A_2(N)$ carries a universal principally polarized abelian
scheme and has a finite surjective forgetful map to the coarse space
$\A_2$; its analytic uniformization and the algebraicity of the
universal family are described in
\cite[Theorem~4.6 and Corollary~4.7]{MilneModuli2013}.
Take an irreducible component of its pullback to $C$ that
dominates $C$, and normalize.  Normalization is finite for a curve over
a field of characteristic zero.  The resulting curve is smooth, and
pullback supplies the abelian scheme.  Removing the preimages of any
finite exceptional set of coarse points leaves a finite surjective map
onto the complementary open of $C$.

All these varieties, components, morphisms, and level structures are
defined over a number field after a fixed enlargement.  The fibers of
a fixed finite map have bounded length; on a dense open this length is
its degree.  The residue field degree of any point in such a fiber is
at most that length.  This proves \eqref{geo:degree-descent}, with the
degrees of any additional fixed covers included in $b$.  It also shows
that the image of every omitted finite set is finite.  An infinite
coarse locus therefore has infinitely many lifts after every such
fixed omission, and we choose one of them as $t_0$.  Later covers may
be supplied with a fixed lift of this point; we then reuse $S,t_0$ for
the refined data.  Open charts and frames used at the center are
chosen to contain it.

A smooth curve over a number field has a smooth projective completion.
The point $t_0$ is an interior point of the moduli family, and its
divisor has positive degree on that completion.  This construction
uses no geometric boundary point and imposes no condition on the
arithmetic reduction of $\mathscr A_{t_0}$.
\end{proof}

Fix a number field $K_C$ defining the original coarse curve.  There
are only finitely many coarse conjugate curves
\begin{equation}\label{geo:coarse-curves}
 C_1,\ldots,C_e\subset\A_2,
\end{equation}
obtained by applying embeddings of $K_C$ into $\Qbar$.  All conjugate
families used below have coarse images in this list.  Enlarging the
field to define centers, levels, or individual endomorphisms does not
add coarse images: the image of $C$ depends only on the restriction
of the embedding to $K_C$.  We always retain the closed coarse curves
in \eqref{geo:coarse-curves}, even when their working parameter curves
have been restricted to open subsets.

\subsection{The five-dimensional orthogonal system}

Let $H=R_1\mathscr A_*\Q$ denote rational covariant first homology.
The polarization gives an alternating form $\psi$, with its usual
Tate twist in the different realizations.  For an operator $D$, write
$D^\dagger$ for the adjoint defined by
$\psi(Dv,w)=\psi(v,D^\dagger w)$.  Define
\begin{equation}\label{geo:operator-space}
 \cP=\{D\in\End(H):D^\dagger=D,\ \Tr_H(D)=0\},
 \qquad \langle D,E\rangle=\Tr_H(DE).
\end{equation}
These definitions commute with all the realization functors used
later.  On actual endomorphisms, $\dagger$ is the Rosati involution
of the given principal polarization.

\begin{proposition}\label{geo:five-space}
The system $\cP$ has rank five and a split nondegenerate trace pairing.
Conjugation induces a surjective morphism
\[
 \GSp_4\longrightarrow\SO(\cP)
\]
with scalar kernel.  Its restriction to $\Sp_4$ is the spin covering,
with kernel $\{1,-1\}$, and $\cP$ is absolutely irreducible.
For every odd prime $\ell$, the lattice
\begin{equation}\label{geo:integral-lattice}
 \Lambda_\ell=
 \{D\in\End_{\Z_\ell}(T_\ell\mathscr A_t):
       D^\dagger=D,\ \Tr(D)=0\}
\end{equation}
is free of rank five with unimodular trace pairing.

If $\mathscr A_t$ is isogenous to the square of a non-CM elliptic
curve, then
\begin{equation}\label{geo:hodge-plane}
 F_t=\End^0(\mathscr A_t)\cap\cP_t
\end{equation}
is a rational two-dimensional plane on which the trace pairing is
positive definite.  Here the intersection denotes Hodge
endomorphisms in the Betti realization, equivalently actual rational
endomorphisms in every realization.
\end{proposition}

\begin{proof}
Use a symplectic basis, with matrix
$J=\left(\begin{smallmatrix}0&I_2\\-I_2&0\end{smallmatrix}\right)$.
The equations $D^{\mathsf t}J=JD$ and $\Tr D=0$ give exactly
\begin{equation}\label{geo:matrix-model}
 D=\begin{pmatrix}A&uK_2\\vK_2&A^{\mathsf t}\end{pmatrix},
 \qquad
 A=\begin{pmatrix}a&b\\c&-a\end{pmatrix},
 \qquad K_2=\begin{pmatrix}0&1\\-1&0\end{pmatrix}.
\end{equation}
Multiplication gives
\begin{equation}\label{geo:trace-form}
 D^2=(a^2+bc-uv)I_4,
 \qquad \Tr(D^2)=4(a^2+bc-uv).
\end{equation}
Thus the form is the orthogonal sum of a nonzero one-dimensional
form and two hyperbolic planes.  Its Gram determinant in the five
coordinates displayed above is $64$, a unit over every $\Z_\ell$
with $\ell$ odd.  A principal polarization makes the Tate module
symplectically unimodular, so the same calculation proves
\eqref{geo:integral-lattice}.

Conjugation by a symplectic similitude preserves both conditions in
\eqref{geo:operator-space} and the trace pairing.  Since $\GSp_4$ is
connected, its orthogonal image has determinant one.  The induced
action of $\Sp_4$ is nontrivial, as is already seen by conjugating the
matrices in \eqref{geo:matrix-model}.  The simple Lie algebra
$\mathfrak{sp}_4$ consequently has zero kernel in this action.  Both
$\Sp_4$ and $\SO_5$ have dimension ten, so the image is all of
$\SO_5$.  The finite kernel is central in the connected group
$\Sp_4$ and is exactly $\{1,-1\}$.  Over an algebraic closure a
similitude is a scalar times a symplectic element, which proves the
scalar-kernel assertion for $\GSp_4$ and identifies the restriction
as the spin covering.  The natural representation of $\SO_5$ is
irreducible over an algebraically closed field of characteristic
zero: the orbit of any nonzero vector, whether isotropic or not,
spans the entire quadratic space.  This proves absolute
irreducibility.  Equivalently, $D\mapsto\psi(D\cdot,\cdot)$
identifies the self-adjoint space with alternating two-forms, and
the trace-free summand with the primitive exterior-square
representation.

For the last assertion, put $A=\mathscr A_t$.  The equivalence between
complex abelian varieties up to isogeny and polarizable rational
Hodge structures of type $(-1,0),(0,-1)$ identifies
$\End^0(A)$ with the Hodge endomorphisms of $H_1(A,\Q)$
\cite[Theorem~4.1]{MilneShimura1995}.  Write
$H_1(A,\Q)=H_1(E,\Q)\otimes\Q^2$ using the given unpolarized
isogeny.  Since $E$ has no CM, its endomorphism algebra is $\Q$;
the Hodge endomorphism algebra of $A$ is therefore $\Mat_2(\Q)$.
An alternating polarization on this tensor product has the form
$\psi_E\otimes B$, where $B$ is a symmetric rational form on
$\Q^2$.  Indeed, each of its four matrix entries is a Hodge
alternating pairing on $H_1(E,\Q)$, hence a scalar multiple of
$\psi_E$; alternation of the total form makes these four scalars
symmetric.  Positivity of the polarization, after choosing the sign
of $\psi_E$, makes $B$ positive definite.  The Rosati involution on
$\Mat_2(\Q)$ is consequently the adjoint for $B$.

Its self-adjoint part has dimension three and contains the identity.
The homology trace on this algebra is twice the ordinary matrix
trace, so imposing trace zero leaves dimension two.  Over $\R$,
conjugation by the positive square root of $B$ turns a self-adjoint
matrix $D$ into a symmetric matrix $D_0$.  Hence, for $D\ne0$,
\[
 \Tr_{H_1(A)}(D^2)=2\tr(D_0^2)>0.
\]
This proves the asserted positivity for every principal
polarization, without requiring the initial isogeny from $E^2$ to
preserve it.
\end{proof}

Fix a labeled pair $\mathbf B=(B_1,B_2)$ of linearly independent
integral endomorphisms spanning $F_0:=F_{t_0}$; clear denominators
in a rational basis to obtain it.  At any target and prime, the
plane lattice used for reduction is
\begin{equation}\label{geo:saturated-plane}
 \Pi_{t,\ell}=(F_t\otimes\Q_\ell)\cap\Lambda_\ell.
\end{equation}
It is saturated and has rank two, so its image in
$\Lambda_\ell/\ell\Lambda_\ell$ has dimension two.  We do not
identify this image with the span of a chosen pair reduced modulo
$\ell$: that pair can fail to be a lattice basis.  At the fixed
center, the restriction of the pairing to
$\Pi_{t_0,\ell}$ is nondegenerate modulo $\ell$ outside finitely
many primes.  To see finiteness, intersect $F_0$ with the integral
Betti operator lattice, choose a $\Z$-basis of this intersection,
and exclude the prime divisors of its nonzero Gram determinant.
Betti--Tate comparison gives precisely the lattices above.

\subsection{Full and joint monodromy}

\begin{proposition}\label{geo:monodromy}
The connected algebraic monodromy of $H$ is $\Sp_4$, and that of
$\cP$ is $\SO_5$.  This remains true on every fixed finite cover,
on a dense open, and for the conjugate families under consideration.
The same connected monodromy is obtained by pullback to a smooth
connected algebraic variety mapping dominantly to the base curve.

Suppose a smooth connected algebraic variety $W$ maps dominantly to
two such curves.  The connected joint monodromy on the two
five-dimensional systems is either $\SO_5\times\SO_5$, or the
coarse joint image of $W$ is contained in a polarized isogeny
correspondence.  If only the first projection varies, its connected
monodromy is still $\SO_5$.
\end{proposition}

\begin{proof}
Hodge genericity in Siegel moduli means that the generic
Mumford--Tate group of $H$ is $\GSp_4$.  The variation is pure and
polarizable and comes from an abelian scheme over a smooth
algebraic curve.  Thus it is among the good variations to which
the connected-monodromy normality theorem applies: connected
algebraic monodromy is a normal subgroup of the derived generic
Mumford--Tate group
\cite[\S5, Theorem~1]{Andre1992}.  Since $\Sp_4$ is almost simple,
the connected monodromy is trivial or all of $\Sp_4$.

Here is a direct exclusion of the trivial possibility.  It would
make the monodromy group finite, and a finite cover would trivialize
the integral local system.  The period map would then be a
single-valued holomorphic map from a smooth algebraic curve to
Siegel space.  Realize Siegel space as a bounded domain in
$\C^3$.  Each coordinate of this map is bounded near each of the
finitely many missing points of the smooth projective completion,
so has a removable singularity there.  Every extended coordinate
is constant on the projective curve.  The family would have
constant moduli image, contrary to its domination of $C$.
This also explains concretely the applicability of
\cite[\S6, Remark~1]{Andre1992}.  The monodromy conclusion is
independent of the Mumford--Tate group at the chosen center.
\Cref{geo:five-space} now gives the assertion on $\cP$.

A nonconstant morphism between smooth algebraic curves restricts,
after removal of finitely many points, to a finite etale morphism
onto a dense open.  Its topological fundamental group therefore
has image of finite index in that of this open; the latter maps
surjectively to the fundamental group of the original curve.
Finite-index subgroups have the same connected Zariski closure.
This proves invariance under finite covers and open restrictions.
For a dominant map from a higher-dimensional smooth variety,
choose an algebraic curve in that variety whose projection is
nonconstant, and normalize it.  Such a curve is obtained by
successive general affine hyperplane sections through a point
where the differential of the projection is nonzero.  Its
monodromy is contained in the pullback monodromy and is already
full by the curve case.  This proves the pullback assertion.

For completeness, full monodromy persists under field conjugation
without a choice of a conjugation-compatible complex period map.
Choose a prime $\ell$.  The Betti monodromy is Zariski dense in
$\Sp_4$, so its image on $H\otimes\Q_\ell$ is Zariski dense in
$\Sp_{4,\Q_\ell}$.  The comparison of topological and geometric
etale fundamental groups identifies the latter image with a dense
subgroup of the geometric etale monodromy image.  A field
conjugation identifies the corresponding geometric etale families
and their symplectic Tate systems.  Applying the same comparison
on the conjugate complex curve proves that its Betti algebraic
monodromy is again full.  In particular its coarse image is Hodge
generic: otherwise its monodromy would lie in the derived group
of a proper generic Mumford--Tate subgroup.  One can also use
the preservation of Siegel special subvarieties by conjugation.

It remains to prove the assertion concerning $W$.  Work at a
generic Mumford--Tate point of the direct sum of the two pulled
back homology variations, in rational symplectic bases.  Both
projections of its connected algebraic monodromy are full, by the
pullback assertion just proved.  Pass to the adjoint groups
$G_1,G_2$, each isomorphic to $\operatorname{PGSp}_4$ and to the
corresponding $\SO_5$.  Let $M\subset G_1\times G_2$ be the
connected joint image.  The kernel of either projection is a
normal subgroup of the other adjoint simple group.  Thus either
one kernel is the whole factor, giving $M=G_1\times G_2$, or both
are trivial and
\begin{equation}\label{geo:graph-group}
 M=\{(x,\phi(x)):x\in G_1\}
\end{equation}
for a $\Q$-isomorphism $\phi:G_1\to G_2$.  This is the group
form of Goursat's argument; no passage to a merely real graph has
been made.

The remaining task is to lift this adjoint graph to a positive rational
similitude identifying the homology Hodge structures.  This requires
both the graph relation and the polarization characters.
Let $T$ be the generic Mumford--Tate group of the direct sum.
The normality theorem applies on the smooth algebraic variety
$W$ and says that $T$ normalizes connected monodromy.  Its
adjoint image therefore normalizes \eqref{geo:graph-group}.
The normalizer of this graph in $G_1\times G_2$ is the graph
itself.  Indeed, normalization by $(u,v)$ says that
$v\phi(x)v^{-1}=\phi(u)\phi(x)\phi(u)^{-1}$ for all $x$;
the trivial center of $G_2$ gives $v=\phi(u)$.

The root diagram of type $C_2$ has no nontrivial automorphism,
so every automorphism of its adjoint split group is inner.
The conjugating element is unique, hence is rational when the
automorphism is rational.  Finally the sequence
\[
 1\longrightarrow\Gm\longrightarrow\GSp_4
   \longrightarrow\operatorname{PGSp}_4\longrightarrow1
\]
and $H^1(\Q,\Gm)=1$ lift it to a rational similitude
$g\in\GSp_4(\Q)$.  Consequently the two Hodge homomorphisms,
after conjugation by $g$, agree in the adjoint group.  Their
remaining difference is a scalar character $c$ of the Deligne
torus.  Each polarization is a morphism of Hodge structures
$\psi_i:H_i\otimes H_i\to\Q(1)$.  Thus the full similitude
character $\nu\circ h_i$ is the character of the same Tate
structure $\Q(1)$ for both $i$, on the whole Deligne torus.
Since a scalar $c$ has multiplier $c^2$, the equality
$h_2=c\,g h_1g^{-1}$ therefore gives $c^2=1$.
The Deligne torus is connected as an
algebraic group, whence $c=1$.  We have obtained a rational Hodge
isomorphism on $H_1$, not just an identification of its adjoint
representation.

Its multiplier is positive.  For if $J_i$ are the complex
structures at this point, then $gJ_1=J_2g$, and positivity of
the two polarization metrics in
\[
 \psi_2(gv,J_2gv)=\nu(g)\psi_1(v,J_1v)
\]
forces $\nu(g)>0$.  By the rational Hodge equivalence
\cite[Theorem~4.1]{MilneShimura1995}, an integral multiple of
$g$ is a polarized isogeny.

This same $g$ describes an analytic open of the joint image.
Indeed, trivialize the rational local system in a simply
connected open.  For each of its countably many rational tensors,
the locus where it is Hodge is closed analytic.  A point outside
the union of those loci that are proper has the generic
Mumford--Tate group.  Every tensor Hodge at that point is
therefore Hodge throughout the open.  By the tensor-stabilizer
description of the Mumford--Tate group, all Hodge homomorphisms
there factor through the transported generic group $T$.
Applying the graph relation
for that group makes $g$ a Hodge isomorphism throughout the
open.  The joint period map is therefore locally in the graph
of $g$, and its coarse image lies locally in the Hecke
correspondence $T_g$.  This correspondence is closed and
algebraic: it is the image of the finite polarized-isogeny
moduli correspondence, or equivalently of the quotient for
$\Gamma\cap g^{-1}\Gamma g$, where
$\Gamma=\Sp_4(\Z)/\{1,-1\}$.  An algebraic equation for $T_g$
vanishing on a nonempty analytic open of the irreducible image
vanishes on that image.  Thus the entire joint image lies in
$T_g$.  The claim when the second projection is constant follows
already from full monodromy of the first projection.
\end{proof}

\subsection{Finitely many dominating correspondence types}

For $g\in\GSp_4(\Q)$ of positive multiplier, write $T_g$ for its
coarse Hecke correspondence in $\A_2\times\A_2$.  Scalar multiples
and integral changes of symplectic markings on either side give
the same correspondence.  A \emph{dominating relation} between
closed curves $C_i,C_j$ means an irreducible algebraic curve
contained in $(C_i\times C_j)\cap T_g$ with both projections
dominant.

\begin{proposition}\label{geo:hecke-list}
For the fixed finite collection \eqref{geo:coarse-curves}, only
finitely many coarse Hecke correspondences admit a dominating
relation between a pair of these curves.  In particular there is
a fixed finite list $\mathcal T$ containing every exceptional
joint image in \Cref{geo:monodromy}.  We may take this list closed
under exchanging the two factors and choose, for every member,
positive integral multipliers for isogenies realizing all its
pairs.  The list and these multipliers are fixed before any
additional auxiliary levels are chosen.
\end{proposition}

\begin{proof}
We prove finiteness, including the analytic stabilizer assertion
on which it depends.  Use one fixed torsion-free principal level $N$ as
in \Cref{geo:setup}.  There are finitely many irreducible closed
curves $D_\alpha\subset\A_2(N)$ lying above the curves $C_i$.
All of them have full connected monodromy by
\Cref{geo:monodromy}.  Write $\mathfrak H_2$ for Siegel space and
\[
 G=\GSp_4(\R)^+/\R^\times
     \simeq\Sp_4(\R)/\{1,-1\},
 \qquad \Gamma_N\subset G
\]
for the effective level group.  The map
$\pi_N:\mathfrak H_2\to\A_2(N)(\C)$ is a covering: the choice
of torsion-free level eliminates stabilizers.

\paragraph{Lift components and their integral stabilizers.}
Remove the finitely many singular points of $D_\alpha$, and
choose a connected component $U_\alpha$ of its inverse image.
Let $Y_\alpha$ be its closure in $\mathfrak H_2$.  It is a
closed irreducible analytic curve.  To check this last assertion,
the inverse image of $D_\alpha$ is locally analytically identical
to $D_\alpha$ under the covering.  It has a locally finite
collection of curve branches, and closure of a chosen connected
smooth part joins precisely its continued branches.  The
normalization of that closure is connected, so the closure is
irreducible.  Conversely every irreducible component arises this
way.  The deck group acts transitively on these components, by
path lifting along the connected nonsingular part of
$D_\alpha$.

Put $\Delta_\alpha=\Gamma_N\cap\Stab_G(Y_\alpha)$.  This is
also the subgroup preserving $U_\alpha$, and covering-space
monodromy identifies it with the image of
$\pi_1(D_\alpha^{\mathrm{reg}})$ in $\Gamma_N$.  In particular
it is Zariski dense in $G$.  Moreover
\begin{equation}\label{geo:lift-quotient}
 \Delta_\alpha\backslash U_\alpha
       \simeq D_\alpha^{\mathrm{reg}}.
\end{equation}
The identification follows because any two lifts of a point in
the chosen component differ by a deck transformation preserving
that component.  Removing additional finite sets, if needed,
has no effect on these statements.

\paragraph{Closedness and discreteness of the real stabilizer.}
Let $L_\alpha=\Stab_G(Y_\alpha)$.  If $g_n\in L_\alpha$ and
$g_n\to g$ in $G$, closedness of $Y_\alpha$ gives
$gY_\alpha\subset Y_\alpha$.  Applying the same argument to
$g_n^{-1}$ proves the reverse inclusion.  Thus $L_\alpha$ is a
closed subgroup of the real Lie group $G$.

Its Lie algebra is invariant under $\Ad(\Delta_\alpha)$.
The stabilizer of a linear subspace of $\mathfrak{sp}_4(\R)$
under the adjoint action is algebraic.  Zariski density therefore
makes $\Lie(L_\alpha)$ invariant under all of $G$, hence an
ideal in the simple real Lie algebra $\mathfrak{sp}_4(\R)$.
It is zero or the entire Lie algebra.  The second possibility
would imply $L_\alpha=G$, since $G$ is connected, and would make
the nonempty curve $Y_\alpha$ invariant under the transitive
action of $G$ on the three-dimensional domain $\mathfrak H_2$.
This is impossible.  Thus $\Lie(L_\alpha)=0$ and $L_\alpha$
is discrete.

\paragraph{Finite area of the integral quotient.}
Equip $\mathfrak H_2$ with its invariant Hermitian metric, and
use the induced area on the regular part of $Y_\alpha$.  We
claim
\begin{equation}\label{geo:finite-area}
 \operatorname{area}(\Delta_\alpha\backslash
                         Y_\alpha^{\mathrm{reg}})<\infty.
\end{equation}
By \eqref{geo:lift-quotient}, it suffices to estimate the induced
metric on the normalization of $D_\alpha$.  Its smooth
projective completion adds finitely many punctures.

Here is the required uniform disk estimate.  Move the value at
zero of a holomorphic disk map into $\mathfrak H_2$ to $iI_2$
by an element of $G$, and use the Cayley transform to the
bounded domain of symmetric two-by-two matrices of operator
norm less than one.  The resulting map $Z$ has $Z(0)=0$.
Scalar Schwarz applied to $u^*Z v$ for unit vectors $u,v$ gives
$\|Z'(0)\|_{\mathrm{op}}\leq1$.  At the origin the invariant
metric is a fixed positive multiple of
$\tr(dZ\,dZ^*)$.  Hence the derivative has a uniform bound.
Precomposing with disk automorphisms gives, with an absolute
constant $c$, the metric inequality
\[
 f^*ds^2_{\mathfrak H_2}
       \leq c\frac{|dz|^2}{(1-|z|^2)^2}
\]
for every holomorphic disk map $f$.

On a punctured disk $0<|q|<R$ in the normalization, lift its
universal covering disk to $\mathfrak H_2$ and apply this
inequality.  It descends to
\begin{equation}\label{geo:cusp-area}
 ds^2\leq c'\frac{|dq|^2}
 {|q|^2\log^2(R/|q|)}.
\end{equation}
Its area integral for $0<|q|<r<R$ is bounded by a constant
times
\[
 \int_0^r\frac{d\rho}{\rho\log^2(R/\rho)}
       =\frac{1}{\log(R/r)}<\infty.
\]
The compact remainder has finite area by smoothness.  Zeros of
the derivative only decrease the area, and singular points and
their discrete lifts have area zero.  This proves
\eqref{geo:finite-area}.

\paragraph{Finite index in the full stabilizer.}
The action of $G$ on $\mathfrak H_2$ is proper: this symmetric
space is $G/K$ with $K$ compact.  Its restriction to the closed
discrete subgroup $L_\alpha$ is consequently proper, with finite
point stabilizers.  Choose a smooth point $y\in Y_\alpha$ and
write $F=(L_\alpha)_y$, a finite group.  Properness provides a
sufficiently small neighborhood $U$ of $y$ in the regular curve
of positive area $a$, such that $hU\cap U=\varnothing$ unless
$h\in F$.

For representatives $h_i$ of distinct double cosets in
$\Delta_\alpha\backslash L_\alpha/F$, the images of $h_iU$
in $\Delta_\alpha\backslash Y_\alpha$ are disjoint.  Indeed,
an intersection would imply
$h_i^{-1}\delta h_jU\cap U\ne\varnothing$ for some
$\delta\in\Delta_\alpha$, hence
$h_i^{-1}\delta h_j\in F$.  Each image has area at least
$a/|F|$, since only members of a conjugate of $F$ can identify
points within it.  Infinitely many such double cosets would
contradict \eqref{geo:finite-area}.  There are therefore finitely
many of them; because $F$ is finite, this proves
\begin{equation}\label{geo:finite-index}
 [L_\alpha:\Delta_\alpha]<\infty.
\end{equation}

\paragraph{Transporters and Hecke double classes.}
For two fixed lift components, consider the real transporter
\[
 \mathcal R_{\alpha\beta}
  =\{g\in G:gY_\alpha=Y_\beta\}.
\]
If it is nonempty and $g_0$ belongs to it, then
$\mathcal R_{\alpha\beta}=g_0L_\alpha$.
By \eqref{geo:finite-index} this set has finitely many right
$\Delta_\alpha$-cosets.  Among the cosets that meet the image
of $\GSp_4(\Q)^+$ choose rational representatives.  These give
finitely many integral double classes for all rational elements
of the transporter, because $\Delta_\alpha\subset\Gamma_N$
and passage to the full integral group only identifies more
classes.

Now suppose $T_g$ admits a dominating relation between two
coarse curves.  Over a smooth point of that relation, both
projections have nonconstant local parametrizations.  Lift them
to the fixed torsion-free level and then locally to Siegel space.  A
local sheet of the Hecke correspondence identifies the two
lifted arcs by a rational positive similitude in the integral
double class of $g$.  The arcs lie on some translates of
$Y_\alpha,Y_\beta$; applying deck transformations makes these
the chosen components.  The translated similitude takes a
nonempty open arc of $Y_\alpha$ into $Y_\beta$.  Its image of
$Y_\alpha$ and $Y_\beta$ are closed irreducible analytic curves
with a one-dimensional intersection, so the analytic identity
principle makes them equal.  The similitude thus belongs to the
transporter just considered.  There are finitely many pairs
$(\alpha,\beta)$ and finitely many double classes for each.
This proves the claimed finiteness of coarse correspondences.

Finally choose a rational similitude representative $g$ of each
retained class.  Multiplying it by a positive integer clears its
matrix denominators.  Its multiplier $m$ is then a positive
integer: applying the equality of alternating forms to a pair
of integral vectors of pairing one shows that the multiplier
is integral.  It yields an isogeny $f$ with
$f^*\lambda'=m\lambda$ at every pair of its correspondence.
This also follows from the finite moduli description by kernels
of such isogenies in the $m$-torsion.  Repeat the choice for the
inverse representatives to include both directions.  There are
only finitely many resulting positive integers.  The construction
has involved only the fixed coarse curves and an initial torsion-free
level; further levels change their parametrizations, not these
curves or the list.  Thus all multipliers are fixed before the
auxiliary incidence primes are selected.
\end{proof}

The finiteness just proved concerns dominating relations.  It places
no bound on the isogeny degrees of isolated pairs of points on the
two curves.  A pair on a retained correspondence, whether or not it
lies on a dominating component, does have an isogeny of that
correspondence's fixed multiplier.  Also, if an irreducible
algebraic joint image is contained in a retained correspondence on
a dense open, its whole image, including any specified target
pair where the maps are defined, is contained there by closedness.
These are precisely the two uses of the finite list in
\Cref{loc:equations,nid:nonidentity}.
